\section{Conclusion}
\label{sec:conclusion}
% \vspace{-10pt}
We introduce \textbf{AgentJudgeBench}, a benchmark measuring LLM-judge reliability on structured, dependency-driven tool-calling, where existing LLM-as-judge work offers little calibration, via three design choices: difficulty-stratified rewrites (easy/medium/hard, preserving the ground-truth trace), per-record DAG-topology annotation across six patterns, and a paired with-GT/without-GT evaluation protocol.

Across $321{,}648$ completed evaluations, four patterns emerge. First, alignment degrades monotonically with difficulty, approximately $1.5\times$ faster without ground truth than with it. Second, hard without-GT alignment converges to a $77$--$82\%$ band regardless of judge capacity, a task-level ceiling confirmed by the C2 recalibrated-prompt ablation (${\leq}{+1.0}$\,pp for capable generators). Third, GT exposure is counterproductive for frontier judges (Gemini-2.5-Pro: $-3.9$\,pp; GPT-5.4: $-1.5$\,pp), consistent with over-anchoring. Fourth, prompt structure is impactful but judge/generator-dependent: on the original (Qwen3-32B, Llama-3.3-70B) pairing, structured rubrics add $+4.8$--$+6.5$\,pp against ${\leq}0.3$\,pp for CoT reasoning, but a second pairing (QwQ-32B, SmolLM3-3B) shows a smaller effect that reverses on hard queries (Limitations); temperature is robustly negligible across both pairings (${\leq}0.25$\,pp spread). The practical guidance is direct: \textbf{with ground truth, use QwQ-32B for closest agreement with the programmatic reference, or GPT-OSS-120B for closest agreement with human judgement (Appendix~\ref{app:prog-judge-validation}); without it, frontier judges lead narrowly, but the without-GT ceiling limits the choice.} Limitations and deployment guidance are in the Limitations section and Appendix~\ref{app:decision-guide}.